\chapter{Data and Reference Decision Framework}
\label{chap:data-reference-framework}

The operational problem addressed in this thesis starts from a stream of
timestamped satellite signal measurements and asks whether the link condition
justifies activating a backup resource. Before any model can be trained, this
problem requires a precise data and decision framework: the raw files must be
converted into a common signal representation, invalid observations and gaps
must be handled consistently, continuous acquisition segments must be defined,
and the reference switch sequence must be constructed from explicit rules.

This chapter defines that framework. Its purpose is not only to prepare data
for modeling, but to fix the operational meaning of the quantities evaluated
later in the thesis. The same data layer determines what counts as an available
observation at prediction time, how heterogeneous sources are reduced to the
signal-only schema, how candidate fade periods are identified, and how model
outputs are compared against a common Perfect Switch reference. The experimental
split and validation protocol then follow from this framework, because model
comparisons are only meaningful when they share the same cleaned data,
timestamp grid, reference decision, and post-processing policy.

\section{Data Preparation Pipeline Overview}
\label{sec:data-preparation-pipeline-overview}

The complete preparation flow is summarized in
\cref{fig:data-preparation-pipeline}. The diagram is intended as a map for the
rest of the chapter: each block is introduced here once, and the following
sections describe the corresponding assumptions and implementation choices in
detail.

\begin{figure}[p]
\centering
\vspace*{-0.35cm}

\resizebox{0.98\textwidth}{!}{%
\begin{tikzpicture}[
    >=Latex,
    common/.style={
        draw,
        rounded corners,
        fill=gray!12,
        align=center,
        text width=4.15cm,
        minimum height=1.05cm,
        inner sep=6pt,
        font=\footnotesize
    },
    prep/.style={
        draw,
        rounded corners,
        fill=yellow!18,
        align=center,
        text width=4.15cm,
        minimum height=1.05cm,
        inner sep=6pt,
        font=\footnotesize
    },
    lane/.style={
        draw,
        rounded corners,
        align=center,
        text width=4.85cm,
        minimum height=1.35cm,
        inner sep=7pt,
        font=\footnotesize
    },
    eval/.style={
        draw,
        rounded corners,
        fill=purple!12,
        align=center,
        text width=4.25cm,
        minimum height=1.05cm,
        inner sep=6pt,
        font=\footnotesize
    },
    arrow/.style={->, thick, gray!75!black},
    ararrow/.style={->, very thick, blue!70!black},
    clparrow/.style={->, very thick, green!50!black},
    lfdarrow/.style={->, very thick, orange!80!black},
    sparrow/.style={->, very thick, purple!75!black},
    note/.style={align=center, font=\footnotesize}
]

% ============================================================
% Introductory note
% ============================================================

\node[
    note,
    text width=12.5cm
] at (5.10,0.65)
{
Common signal-only preparation: all tasks start from the same cleaned
univariate signal representation, and no future observations are used
as model inputs.
};

% ============================================================
% Common preprocessing
% ============================================================

\node[common] (raw) at (0.00,-0.85)
    {Raw satellite signal files};

\node[common] (load) at (5.10,-0.85)
    {Signal-only loading: \texttt{Time}, \texttt{Signal}};

\node[common] (clean) at (10.20,-0.85)
    {Timestamp parsing, sorting, cleaning, and schema validation};

% Slightly raised preprocessing row
\node[prep] (segments) at (10.20,-2.70)
    {Sampling-time checks, conservative gap handling,
     and acquisition segments};

\node[prep] (events) at (5.10,-2.70)
    {Candidate fade events, event windows, and quality filtering};

% Slightly lowered to create more separation
\node[
    prep,
    text width=5.65cm
] (split) at (5.10,-4.45)
    {Valid sample index, split arrays, metadata,
     and external-holdout split};

% Common pipeline arrows

\draw[arrow] (raw) -- (load);
\draw[arrow] (load) -- (clean);
\draw[arrow] (clean.south) -- (segments.north);
\draw[arrow] (segments.west) -- (events.east);
\draw[arrow] (events.south) -- (split.north);

% ============================================================
% Task branching
% ============================================================

% Slightly lowered together with split
\coordinate (branch) at (5.10,-5.50);

\draw[arrow] (split.south) -- (branch);

% ============================================================
% Task-specific blocks
% ============================================================

\node[
    lane,
    fill=blue!10
] (ar) at (1.80,-6.85)
{
\textbf{Autoregressive forecasting}\\[3pt]
Context \(\rightarrow\) trajectory
};

\node[
    lane,
    fill=green!12
] (clp) at (8.40,-6.85)
{
\textbf{Current-level persistence}\\[3pt]
Context \(\rightarrow\) duration
};

\node[
    lane,
    fill=orange!14
] (lfd) at (1.80,-9.15)
{
\textbf{Long-fade detection}\\[3pt]
Event context \(\rightarrow\) class
};

\node[
    lane,
    fill=purple!10
] (sp) at (8.40,-9.15)
{
\textbf{Survival persistence}\\[3pt]
Context \(\rightarrow\) survival curve
};

% ============================================================
% Central branch
% ============================================================

\coordinate (branchbottom) at (5.10,-11.00);

\draw[
    thick,
    gray!75!black
] (branch) -- (branchbottom);

% Horizontal task arrows

\draw[ararrow]
    (5.10,-6.85) -- (ar.east);

\draw[clparrow]
    (5.10,-6.85) -- (clp.west);

\draw[lfdarrow]
    (5.10,-9.15) -- (lfd.east);

\draw[sparrow]
    (5.10,-9.15) -- (sp.west);

% ============================================================
% Shared reference / evaluation
% ============================================================

\node[eval] (derived) at (1.80,-12.25)
{
Task-specific\\
model outputs
};

\node[eval] (perfect) at (8.40,-12.25)
{
Perfect Switch reference\\
from observed signal
};

\node[eval] (post) at (1.80,-13.85)
{
Shared switch\\
post-processing
};

\node[eval] (metrics) at (8.40,-13.85)
{
Switch metrics and\\
cross-task comparison
};

% ============================================================
% Evaluation arrows
% ============================================================

\draw[arrow, rounded corners=4pt]
    (branchbottom)
    -- ++(0,-0.35)
    -| (derived.north);

\draw[arrow]
    (derived.east) -- (perfect.west);

\draw[arrow]
    (derived.south) -- (post.north);

\draw[arrow]
    (perfect.south) -- (metrics.north);

\draw[arrow]
    (post.east) -- (metrics.west);

\end{tikzpicture}%
}

\caption{Overview of the signal-only data preparation and reference-decision
pipeline. Gray and yellow blocks are shared preprocessing steps; colored
blocks are task-specific dataset targets and downstream switch-evaluation
inputs.}

\label{fig:data-preparation-pipeline}
\end{figure}


The key point is that the pipeline is shared up to the point where each
modeling task requires a different target. Autoregressive forecasting needs
past-future signal windows, current-level persistence needs residual duration
relative to the current signal level, long-fade detection needs a grouped-event
class label, and survival persistence needs observed or right-censored
remaining event durations. The raw data cleaning, segment boundaries, external
holdout policy, and switch reference construction remain common, so differences
in the results can be attributed to the task formulation and model family
rather than to incompatible preprocessing.

\section{Available Signal Datasets}

The project uses a collection of satellite fade and beacon signal files
associated with the Fucino ground segment in Italy and with GX4/GX5 satellite
links \cite{telespazio_fucino,satbeams_gx4,satbeams_gx5}. The files were made
available for this thesis as project data rather than as a public benchmark
dataset. They cover different sources and acquisition periods. Each file
contains a sequence of timestamped observations of the link signal or
attenuation level, but the raw format is not homogeneous across sources. Some
files include explicit column headers, while others are plain delimited files
without column names. Some files contain only timestamp and signal columns,
whereas a small subset includes additional environmental columns. The first
requirement of the data layer is therefore to map these heterogeneous files
into a common representation before any task-specific dataset is built.

The raw inventory used in the final experiments contains eight loaded signal
datasets. Their time coverage ranges from several weeks to several months, with
the earliest acquisition starting in July 2020 and the latest ending in July
2021. \Cref{tab:raw-dataset-inventory} summarizes the loaded files at the level
needed for the experimental protocol.

\begin{table}[htbp]
  \centering
  \caption{Raw signal datasets used in the project inventory. Row counts refer
  to valid rows after the common loading and validation stage.}
  \label{tab:raw-dataset-inventory}
  \begin{tabular}{lrr}
    \toprule
    Dataset & Valid rows & Duration (hours) \\
    \midrule
    \texttt{uplink\_beacon\_data\_fucino\_GX4.csv} & 481715 & 4081.5 \\
    \texttt{dati\_meteo\_beacon.csv} & 484031 & 4081.5 \\
    \texttt{fc-uplink-fade.csv} & 378191 & 3218.0 \\
    \texttt{fc-downlink-fade.csv} & 380566 & 3215.8 \\
    \texttt{UPLINK-FADE\_Fucino\_GX4\_AprJun2021.csv} & 259992 & 2171.7 \\
    \texttt{UPLINK-FADE\_Fucino\_GX5\_AprJun2021.csv} & 256519 & 2161.0 \\
    \texttt{DOWNLINK-FADE\_Fucino\_GX4\_MarJun2021.csv} & 338281 & 2826.0 \\
    \texttt{DOWNLINK-FADE\_Fucino\_GX5\_AprJun2021.csv} & 256519 & 2161.0 \\
    \bottomrule
  \end{tabular}
\end{table}

All datasets used in this work are associated with measurements from the
Fucino ground segment in Italy. Fucino is one of Telespazio's main space
centres and teleports \cite{telespazio_fucino}. The satellite identifiers in
the file names refer to GX4 and GX5, corresponding to the Inmarsat Global
Xpress spacecraft Inmarsat GX4, catalogued as NORAD 42698, and Inmarsat GX5,
catalogued as NORAD 44801 \cite{satbeams_gx4,satbeams_gx5}. Viasat completed
its acquisition of Inmarsat in 2023, so these spacecraft belong to the current
Viasat/Inmarsat provider context \cite{viasat_inmarsat_acquisition}. The
historical Inmarsat GX names are nevertheless retained in this thesis because
they are the identifiers used in the raw files and in public satellite
catalogues.

Within this context, several files are also related by link direction. In the
dataset names, \emph{uplink} denotes the ground-to-satellite propagation path,
while \emph{downlink} denotes the satellite-to-ground propagation path. Thus,
uplink and downlink files can be interpreted as physically paired measurements
or derived fade series for opposite directions of the same broader
communication link, especially when they refer to the same site and payload
identifier. In this thesis, however, these files are not combined into a
multivariate paired representation. Each file is treated as an independent
univariate signal series, because the final modeling policy requires the same
single input variable, \texttt{Signal}, to be available for every dataset.

The datasets differ not only in their acquisition period, but also in their
maximum observed signal excursions and in the amount of invalid or missing
source rows. They also differ in the availability of auxiliary variables.
\Cref{tab:ignored-source-columns} reports the files whose raw format contains
more than the timestamp and signal fields used by the models.

\begin{table}[htbp]
  \centering
  \caption{Additional source columns found in the raw files and ignored by the
  signal-only data policy.}
  \label{tab:ignored-source-columns}
  \begin{tabular}{lp{8.5cm}}
    \toprule
    Dataset & Ignored source columns \\
    \midrule
    uplink beacon GX4
      & \texttt{Temperature}, \texttt{Pressure}, \texttt{Humidity} \\
    meteo beacon
      & \texttt{airt\_bt}, \texttt{pres\_bt}, \texttt{relh\_bt}, \texttt{temp\_bt} \\
    fc-uplink fade
      & third trailing column, empty in the inspected rows \\
    fc-downlink fade
      & third trailing column, empty in the inspected rows \\
    \bottomrule
  \end{tabular}
\end{table}

The environmental columns would be potentially useful in a broader
multivariate nowcasting study, since atmospheric conditions can affect
attenuation dynamics. They are not used here because they are not available in a
consistent form across all datasets. Including them would force the experiments
to either discard most datasets or compare models trained with different input
information. By contrast, several of the Fucino 2021 files contain exactly two
columns and are loaded through positional column selection. This variability
motivates the strict loading policy described in the next section: downstream
models should not depend on source-specific formatting decisions.

The final experimental protocol treats \texttt{fc-uplink-fade.csv} as an
external holdout. This file is included in the project inventory and in the
final processed datasets, but it is reserved for test evaluation. All remaining
datasets form the development pool used for training and validation. Exact
supervised sample counts are reported with the experimental setup, because they
depend on the task-specific target and validity rules.

\section{Signal-Only Policy}

All models in this thesis follow a strict signal-only policy. After loading and
validation, every dataset is standardized to exactly two columns:

\begin{center}
\texttt{Time}, \texttt{Signal}.
\end{center}

The \texttt{Time} column identifies when the observation was acquired. It is
used for chronological ordering, gap detection, elapsed-duration computation,
and splitting into train, validation, and test sets. It is not used directly as
a model input.
The \texttt{Signal} column is the only measured variable made available to the
models. In the raw files, this quantity represents the signal or fade level
used to describe the link degradation. The exact source formatting may differ,
but after the common loading step all downstream tasks operate on the same
semantic object: a univariate time series of signal values.

This policy is deliberately restrictive. Some raw files contain environmental
measurements, such as temperature, pressure, or humidity, and these variables
could plausibly improve attenuation nowcasting in a dedicated multivariate
study. They are excluded here because they are not available consistently
across the full dataset collection. If such variables were used, the
experimental comparison would no longer isolate the effect of the modeling
formulation. Some models would exploit richer inputs, while others would be
trained on signal-only data or on a reduced subset of files. The resulting
differences could then reflect input availability rather than modeling
quality.

The signal-only policy also reduces the risk of uncontrolled leakage. Metadata
columns, dataset identifiers, event labels, final event durations, and future
event properties are never passed as model inputs. The timestamp is not used as
a standalone calendar covariate: it is not converted into hour-of-day,
day-of-year, or month features. This avoids letting the models exploit
dataset-specific acquisition periods or seasonal correlations that may not
transfer to a new deployment period.

The derived predictors allowed by the final datasets must be observable at the
decision time. Most of them are deterministic functions of the past and present
signal window, such as a recent mean, a local range, a recent slope, or a value
expressed relative to the current signal or to the fade threshold. Some
task-specific datasets may also include online elapsed-time quantities defined
by a threshold state machine, such as time since the first crossing of the
current grouped event. These features still obey the signal-only policy because
the state machine is driven only by past and current \texttt{Signal} values and
their timestamps. They do not include calendar predictors, source-specific
metadata, future event ends, final event durations, or labels.

In practical terms, this means that all task formulations are evaluated under
the same information constraint. Autoregressive forecasting, current-level
persistence, long-fade detection, and survival persistence ask different
questions about the same observed signal history, but none of them receives
additional external predictors. This makes the experimental comparison less
ambitious than a full operational system with meteorological inputs, but it
makes the methodological comparison cleaner and reproducible across all
available files.

\section{Parsing, Cleaning, and Validation}

The first processing stage converts each raw file into a validated
\texttt{Time}/\texttt{Signal} table. This stage is intentionally separated
from all modeling scripts. Its purpose is to make every subsequent task start
from the same chronological signal representation, independently of the
formatting details of the original file.

Raw files may or may not contain column headers. When headers are available,
the configured timestamp and signal columns are selected by name when possible.
When headers are absent, or when positional loading is explicitly required, the
first column is interpreted as \texttt{Time} and the second column as
\texttt{Signal}. If the file contains fewer than two columns, loading fails
with an explicit error. If the file contains more than two columns, the extra
columns are ignored and recorded in the loading metadata.

The parser performs only strict conversion. Timestamps are parsed into datetime
values and signal entries are converted to numeric values. Rows with missing or
unparseable timestamps are excluded, because they cannot be placed in the time
series. Rows with missing, non-numeric, or non-finite signal values are also
excluded, because they cannot be used as observed signal measurements. These
rows are not silently hidden: the corresponding counts are saved for each
dataset and used in the inventory summaries.

Some raw files may also contain suspicious finite values that are technically
numeric but likely represent default missing-data codes. In this project,
signal values at or below -10 are reported as suspected sentinel values. They
are not removed automatically at the parsing stage, because doing so would
embed a domain assumption directly into the generic loader. Instead, they are
kept as observations and made visible through the data-quality metadata. This
keeps the cleaning policy auditable and avoids silently changing the empirical
signal distribution.

After parsing, each dataset is cleaned and temporally segmented. The common
cleaning stage applies the following operations:

\begin{itemize}
  \item sort each dataset chronologically;
  \item resolve duplicate timestamps according to the configured strategy;
  \item assign a deterministic \texttt{dataset\_id} and \texttt{row\_id};
  \item compute the elapsed time from the previous observation;
  \item assign a continuous acquisition \texttt{segment\_id};
  \item save the cleaned signal table and the loading metadata.
\end{itemize}

Duplicate timestamps are handled before segment construction. In the canonical
configuration, duplicate observations with the same timestamp are collapsed by
averaging their signal values. Alternative strategies, such as keeping the
first value, keeping the last value, or raising an error, are supported by the
preprocessing layer, but the final experiments use one configured policy
consistently. This is important because duplicate handling can otherwise
change threshold crossings and event boundaries.

The \texttt{segment\_id} marks continuous acquisition stretches. It is created
from the chronological time differences: whenever the elapsed time between two
successive observations exceeds the configured gap threshold, a new segment is
started. At this stage no missing samples are inserted and no signal values are
interpolated. The cleaned tables therefore remain faithful to the observed
timestamps, while still making acquisition breaks explicit for later
no-leakage checks.

Validation is performed through explicit metadata rather than implicit
assumptions. For every raw file, the pipeline records the number of raw rows,
valid rows, invalid timestamp rows, invalid signal rows, suspected sentinel
rows, duplicate timestamps, ignored columns, and whether positional fallback
was used. These diagnostics make it possible to identify datasets that require
special attention before they influence model training or evaluation.

The output of this stage is the common clean signal layer stored under the
intermediate data directory. Downstream task builders reuse this layer instead
of reading raw CSV files independently. This design prevents each task from
implementing its own parser, duplicate policy, or timestamp convention. As a
result, differences between autoregressive forecasting, current-level
persistence, long-fade detection, and survival persistence are due to the task
definition and model family, not to inconsistent low-level data handling.

\section{Sampling Time, Gaps, and Conservative Imputation}

The final experiments assume a nominal sampling interval
\(\Delta t = 30\) seconds. This interval is used to interpret context lengths,
forecast horizons, switch-time constraints, and duration thresholds. For
example, a context length of \(L=30\) corresponds to approximately 15 minutes
of past observations, while a prediction horizon of \(h=10\) corresponds to
approximately 5 minutes into the future. These conversions are meaningful only
if the sequence is approximately regular, so the treatment of missing
timestamps is part of the experimental methodology rather than a minor
implementation detail.

Let an observed signal segment be represented as an ordered sequence
\[
  (t_1, s_1), (t_2, s_2), \ldots, (t_n, s_n),
\]
where \(t_i\) is the timestamp and \(s_i\) is the corresponding signal value.
For two consecutive observations in the same acquisition segment, define
\[
  d_i = t_i - t_{i-1}.
\]
If \(d_i \leq 1.5\Delta t\), the pair is considered close enough to the
expected sampling rate and no missing sample is inferred. If \(d_i >
1.5\Delta t\), the number of missing samples is estimated as
\[
  m_i = \mathrm{round}\left(\frac{d_i}{\Delta t}\right) - 1.
\]
This rule identifies gaps that are longer than the nominal interval but still
compatible with a small number of missing 30-second samples.

The conservative imputation policy used in the final datasets is:

\begin{itemize}
  \item expected sampling interval: 30 seconds;
  \item maximum gap eligible for imputation: 90 seconds;
  \item maximum number of inserted samples in one gap: two;
  \item interpolation method: linear interpolation in time;
  \item allowed scope: only inside the same continuous acquisition segment.
\end{itemize}

Therefore, a gap is imputed only if \(m_i > 0\), \(m_i \leq 2\), and
\(d_i \leq 90\) seconds. For each missing index \(k=1,\ldots,m_i\), the
inserted timestamp and signal value are
\[
  \tilde{t}_{i,k} = t_{i-1} + k\Delta t,
\]
\[
  \tilde{s}_{i,k}
  =
  s_{i-1}
  +
  \frac{k}{m_i+1}
  (s_i - s_{i-1}).
\]
The reconstructed value is thus a straight-line interpolation between the
previous and next observed samples. No model is used to infer the missing
signal. The method does not attempt to reconstruct hidden attenuation peaks,
fast rain-cell dynamics, or any physical process during the missing interval.
It only repairs very short holes that are consistent with one or two omitted
regular samples.

The most important constraint is that imputation is never allowed across
different acquisition segments. Segment boundaries are defined before
imputation from gaps exceeding the configured continuity threshold. Once two
observations belong to different \texttt{segment\_id} values, no interpolation
can connect them. A later observation after such a break is therefore not
treated as evidence that the signal stayed above or below any threshold during
the unobserved period.

This has direct consequences for target construction. In autoregressive
forecasting, a context-target pair is valid only if the required past and
future observations can be represented on the prepared event-window grid. In
current-level persistence, a recovery observed only after an unbridgeable gap
cannot be used as the true recovery time for a sample before the gap. In
survival persistence, if the end of the fade event is not observed before the
continuous segment ends, the corresponding sample is right-censored rather
than assigned an artificial observed residual duration. The policy therefore
prevents missing data from being converted into false certainty about event
persistence or recovery.

There are differences across tasks, but not across models within the same
task. The imputation policy is applied before model training or evaluation, so
all models belonging to the same task consume the same prepared dataset. GRU,
PatchTST, autoregressive XGBoost, and Chronos are evaluated on the same
autoregressive windows for a given context length and horizon. The
current-level shapelet models and current-level XGBoost share the same
current-level dataset. The long-fade classifiers share the same grouped-event
classification dataset, and the survival models share the same
survival-persistence dataset with observed and censored targets.

The task-level usage is as follows. The autoregressive preparation first
constructs event-centered windows and stores a prepared signal version with
small recoverable gaps filled. This prepared signal is also used by task
builders that start from the shared event-window artifacts, such as
current-level persistence and long-fade detection. Duration-oriented builders
also apply the same small-gap rule to the full cleaned signal inside each
continuous segment when they need to search for a future recovery or event end.
Survival persistence applies the same rule before running its event state
machine, so that very short missing intervals do not artificially fragment an
otherwise continuous fade episode. In every case, imputation remains local,
auditable, and bounded by the same numerical constraints.

Event-window quality is derived from this gap analysis. Windows with acceptable
coverage and only recoverable small gaps can be marked as usable or
warning-level, depending on the number and length of the gaps. Windows with
insufficient coverage, excessive missing runs, or gaps that cannot be repaired
are excluded from the final supervised window index. The retained datasets also
store imputation summaries, including how many rows were inserted, how many
gaps were filled, and the largest imputed gap per group.

The policy is intentionally conservative. It may discard some potentially
informative samples, but it avoids a more serious failure: constructing inputs,
targets, or switch references from time intervals where the signal was not
actually observed. This trade-off is consistent with the operational objective
of the thesis, where reliable switch behavior is more important than
maximizing the number of training windows at any cost.

\section{Candidate Events and Acquisition Segments}

The modeling tasks do not sample uniformly from the complete raw files. Most
timestamps correspond to normal or weakly attenuated propagation conditions,
whereas the operational decision is relevant around fade episodes. The common
preparation stage therefore introduces two related, but distinct, concepts:
\emph{acquisition segments} and \emph{candidate fade events}.

An acquisition segment is a continuous portion of one loaded signal file. After
the signal has been sorted by time, a new segment is started whenever the time
distance between two consecutive valid observations is larger than the maximum
allowed continuity gap. In the canonical configuration this limit is
\(90\) seconds. Formally, if \(d_i=t_i-t_{i-1}\), then observations \(i-1\) and
\(i\) belong to the same segment only when \(d_i \leq 90\) seconds. The resulting
\texttt{segment\_id} is used as a guardrail throughout the pipeline: input
contexts must not cross segment boundaries, and future targets must not be
resolved by jumping over an unobserved acquisition break.

Candidate fade events are instead a practical sampling device. For each dataset,
the preparation code first identifies threshold-crossing timestamps,

\[
  \mathcal{C}_{\theta} = \{t_i : S(t_i) \geq \theta\},
\]

with \(\theta=10.0\) dB in the main experiments. The condition is deliberately
non-strict, so a sample exactly equal to \(10.0\) dB is treated as belonging to
the threshold region. Crossings are then grouped chronologically. A group starts
at its first crossing time \(t_0\), and subsequent crossings remain in the same
candidate event while they occur within the configured grouping horizon from
\(t_0\). The default horizon is three hours. A later crossing outside that
horizon starts a new candidate event.

Each candidate event is anchored at the first threshold crossing, called the
\texttt{event\_timestamp}, and is expanded into an event-centered analysis window,

\[
  [t_0 - 1.5\text{h},\; t_0 + 1.5\text{h}],
\]

using the canonical configuration. This window is not intended to say that the
fade lasts three hours. It is only a local region that contains pre-event
context, the threshold-crossing period, and post-event recovery or follow-up
samples when they are observed. The window may therefore contain both
above-threshold and below-threshold samples.

These candidate events should not be confused with the final reference switch,
nor with the task-specific event definitions used later. They are used to make
the raw time series operationally manageable, to compute coverage and gap
quality checks, and to decide which local portions of the files can generate
supervised samples. Window quality is assessed after bounded imputation: the
pipeline records coverage, maximum gap length, number of inferred missing
samples, number of imputed samples, duplicate timestamps, missing values, and
the number of points above and below the threshold. Events with insufficient
coverage or unrecoverable gaps are excluded from the final window index, while
warning-level events can be retained when this is explicitly enabled.

Task-specific datasets are then built on top of this common signal layer:

\begin{itemize}
  \item autoregressive forecasting uses event-window indices to build past
  context and future trajectory pairs. A training sample is valid only when the
  full context and the full prediction horizon are available on the prepared
  time grid;
  \item current-level persistence searches the full cleaned signal for the
  first future recovery below a level relative to the current signal. The
  candidate window can determine which timestamps are considered, but the target
  duration is not truncated at the event-window boundary;
  \item long-fade detection groups threshold crossings within a configured
  temporal horizon and assigns one binary label to all timestamps in the grouped
  fade episode. Temporary below-threshold oscillations inside the grouped event
  do not split the event;
  \item survival persistence uses a state machine with activation, candidate
  recovery, stable recovery, and right-censoring at segment boundaries. Its
  target is the remaining duration until the end of the whole fade event, not
  the duration of the current signal level.
\end{itemize}

The separation between acquisition segments, candidate events, and task-specific
labels is important. Acquisition segments describe what was continuously
observed. Candidate events describe where the pipeline should look for
operationally relevant samples. The supervised target definition is then chosen
by the task. This design preserves a common data foundation while allowing each
model family to answer a different operational question without silently
changing the meaning of an event.

\section{External Holdout Split}

The final evaluation uses an external-holdout protocol. One complete raw file,
\texttt{fc-uplink-fade.csv}, is reserved for test evaluation. It is not used for
model fitting, hyperparameter selection, early stopping, architecture selection,
normalization-statistic fitting, or decision-rule calibration. All remaining
datasets form the development pool.

The split is therefore performed at two levels. At the file level, the held-out
dataset is assigned entirely to the test split:

\[
  \mathcal{D}_{\text{test}} = \{\texttt{fc-uplink-fade.csv}\}.
\]

The development set is the complement,

\[
  \mathcal{D}_{\text{dev}} = \mathcal{D}\setminus\mathcal{D}_{\text{test}}.
\]

At the event level, each development dataset is split chronologically into
training and validation events. Events are first ordered by their earliest
available input timestamp. The first \(85\%\) of development events are assigned
to training and the remaining \(15\%\) to validation, subject to a minimum-event
guard. If a development dataset contains too few events to support a meaningful
validation subdivision, all its events are kept in training rather than creating
an unstable validation split. In the canonical autoregressive and survival
dataset builders the minimum is three events; in the long-fade configuration it
is two events.

This event-aware split prevents windows from the same physical fade episode from
appearing in both training and validation. It also prevents the external test
file from leaking into the development process. The validation split is used for
model selection: selecting grid-search configurations, choosing early-stopping
epochs, comparing candidate architectures, and calibrating only those decision
rules that are explicitly treated as validation-calibrated. The test split is
evaluated after this selection step. If a test-set sweep is ever performed for
diagnostic purposes, it must be reported as exploratory and not as a
validation-calibrated operating point.

The same principle is applied across tasks, although the exact number of samples
differs because each task defines a different supervised target and therefore a
different set of valid decision timestamps. Autoregressive forecasting produces
past-context and future-trajectory windows. Current-level persistence produces
duration-regression samples. Long-fade detection produces grouped-event
classification samples. Survival persistence produces observed and right-censored
residual-duration samples. In all cases, however, the final test data come from
the same held-out acquisition file.

This protocol is stricter than a random sample split. A random split would place
nearby, highly correlated timestamps on both sides of the split and would inflate
performance estimates. The external-holdout protocol is closer to the intended
use case: a model or decision rule is developed on existing acquisition files and
then applied to a separate file. It should still be interpreted as file-level
generalization within the available Fucino data, not as proof of generalization
to all teleports, satellite links, seasons, or operating conditions.

\section{Perfect Switch Reference}

All model-derived switch decisions are evaluated against a model-independent
reference called the Perfect Switch. The Perfect Switch is not a trained model.
It is a reference decision derived from the observed signal and the configured
operational threshold. It is an oracle in the evaluation sense: it can use the
future observed signal to decide whether a switch should have been active, but
it is never available to a real online model.

The reference starts from an instantaneous outage mask. In the canonical
experiments the threshold is \(\theta=10.0\) dB and comparisons are inclusive:

\[
  M_i = \mathbb{1}\{S(t_i) \geq \theta\}.
\]

This binary mask alone is not the Perfect Switch. It only states whether the
currently observed signal is in the threshold region. The Perfect Switch also
incorporates the operational persistence requirement: a switch should not be
activated for isolated threshold crossings that disappear immediately. With
30-second sampling and \(\texttt{switch\_time}=10\), the reference requires
eleven consecutive above-threshold samples, corresponding to ten 30-second
intervals, i.e. approximately five minutes of persistence.

Equivalently, for each contiguous run of \(M_i=1\), the raw Perfect Switch is set
to one only if the run contains at least
\(\texttt{switch\_time}+1\) samples. Shorter above-threshold runs remain zero.
This produces the stored column \texttt{perfect\_switch\_raw}. A legacy
right-extension step is then applied to enforce the minimum switch-time behavior
used in the previous implementation. Finally, if the switch is already active
and the observed signal is still above threshold, the switch is held active even
if the intermediate rule would otherwise return zero. This final held signal is
the column used as the reference in the experiments:

\[
  \texttt{perfect\_switch} \equiv \texttt{perfect\_switch\_min\_time}.
\]

The pipeline also stores diagnostic columns for auditing this transition:

\begin{itemize}
  \item \texttt{outage\_mask}, the instantaneous threshold mask;
  \item \texttt{perfect\_switch\_min\_island\_old}, the legacy right-extension
  intermediate signal;
  \item \texttt{perfect\_switch\_adjusted}, a separate adjusted diagnostic
  variant.
\end{itemize}

These columns help inspect the transition from threshold mask to reference
switch, but the main switch metrics compare model decisions against
\texttt{perfect\_switch}.

The Perfect Switch is task-scoped in terms of timestamp grid. Each task can
produce decisions only at its native decision timestamps, so the reference is
aligned to those timestamps before computing task-level switch metrics. The
underlying threshold convention and post-processing rule, however, are shared
across tasks. This is why any change to the Perfect Switch definition is a
methodological change: it affects every reported switch metric.

\section{Shared Switch Post-Processing}
\label{sec:shared-switch-post-processing}

The reference framework operates after a task-specific prediction has already
been converted into a raw binary decision. The predictive quantity may be a
forecast trajectory, a duration estimate, a class probability, or a survival
curve, but that distinction belongs to the task formulation. At this stage, all
models expose the same intermediate decision column:

\begin{center}
\texttt{model\_switch\_raw}.
\end{center}

The mapping from a native model output to \texttt{model\_switch\_raw} is
task-specific and is defined with the individual formulations in
\cref{chap:task-formulations}. The data-and-reference layer does not choose
those task thresholds or probability operating points; it only receives the raw
decision sequence and applies the same temporal policy to all methods.

After this task-specific conversion, all methods use the same temporal
post-processing. This is intentionally separated from model training: changing
it changes the evaluation decision rule, but it does not require retraining the
underlying forecasting, regression, classification, or survival model.

The first post-processing step enforces a minimum active switch duration on the
raw decision sequence. Each contiguous island of ones is inspected. If the
island is shorter than \(\texttt{switch\_time}\), it is extended to the right
until it reaches that length, or until the available event grid ends. This rule
does not fill arbitrary zero gaps between two switch islands. It can only make
two islands touch if the right-extension of the first island reaches the second
one.

The second step is stateful and uses only information available at the current
decision time. If the post-processed switch was active at the previous timestamp
and the currently observed signal is still above the operational threshold, the
switch is kept active even if the raw model decision has temporarily returned to
zero. This avoids an unrealistic behavior in which the system switches off
while the fade is still ongoing. Once the observed signal leaves the threshold
region, the switch is no longer forced by this holding rule.

The post-processed decision produced by these two steps is the one used for the
main reported switch metrics:

\begin{center}
\texttt{model\_switch}.
\end{center}

This separation prevents a common source of ambiguity. Parameters used inside a
task-specific conversion rule, such as forecast-horizon requirements or
probability thresholds, are different from \texttt{switch\_time}, which controls
only the temporal post-processing of an already binary decision sequence.

Raw and intermediate switch columns are still saved for diagnostics. For model
outputs, \texttt{model\_switch\_raw} is the direct task-specific binary
decision before post-processing; \texttt{model\_switch\_min\_island\_old} is
the right-extended minimum-duration version; \texttt{model\_switch\_min\_time}
is the held version after the above-threshold rule; and
\texttt{model\_switch\_adjusted} is a separate legacy diagnostic variant. These
columns are useful for understanding whether an error comes from the model
output itself or from the operational post-processing rule. Unless explicitly
stated otherwise, the reported switch metrics use the held minimum-time version,
saved as both \texttt{model\_switch\_min\_time} and \texttt{model\_switch}.

\section{Evaluation Alignment}

The data and reference layer provides the common objects needed for operational
evaluation: the observed signal, the instantaneous outage mask, the
task-aligned Perfect Switch, and the intermediate switch columns generated by
the shared post-processing rule. It does not define the native losses of the
models, select hyperparameters, or choose task-specific probability and
duration thresholds. Those experimental choices are specified in
\cref{chap:task-formulations,chap:experimental-setup}.

For each task, model outputs are first converted into
\texttt{model\_switch\_raw} according to the task formulation. The shared
post-processing then produces \texttt{model\_switch}, which is compared with the
aligned \texttt{perfect\_switch}. The same tables also retain diagnostic fields
for the outage mask, raw and adjusted Perfect Switch variants, and the raw model
decision. These fields make it possible to audit whether a disagreement comes
from the predictive model, from the task-specific conversion rule, or from the
operational post-processing layer.

Metric formulas and comparison denominators are intentionally left to the
experimental setup chapter. The role of the present chapter is narrower: it
defines the signal representation and the reference decision objects that all
tasks must use. This keeps the thesis definitions separated: \cref{chap:background}
introduces general concepts, this chapter defines the data and reference layer,
\cref{chap:task-formulations} defines task-specific model outputs and raw switch
rules, and \cref{chap:experimental-setup} defines how the resulting runs are
selected and compared.
